%----------
%	INTRODUCTION
%----------	

\color{red}[mundo globalizado -> uso de redes sociales -> anonimato -> hate speech -> mujeres]\color{black} \newline

In this time and era, the society of the 21st century has evolved from a manual and analog world to an automatic and digital network where everyone and everything is connected. The times of buying the newspaper or sitting in front of the television to watch the news are decreasing more and more, due to the widespread of social media in the last decade. We live in the era of the sporadic content: tweets, Instagram posts, Tik-Tok videos, where a vast amount of content is being published in a daily basis. This shift has been fueled partly by how easy and anonymous these platforms are to use, and partly by people's desire to share or defend their views. Hence, the social debate has changed from the Agora, churches, town squares and taverns, with a limitation in the capacity and attendees, to social platforms in which every single person with an account can participate in.


\color{red}[incluir anonymity]\color{black} \newline
“online disinhibition effect.” \cite{fox-2015}  online anonymity is associated with increased online and offline hostility.


This is the so called by Guy Debord as \textit{La société du spectacle} (The Society of the Spectacle): "All that once was directly lived has become mere representation". We are constantly being overwhelmed with tons of information, news and content that in the majority of the time are useless, fake or even hate towards a certain person or group. In particular, in the context of mass-media and social networks, human relations and hence, social debate are becoming less authentic and personal. This is due to the fact that 


\color{red}[incluir algo sobre hate speech a mujeres]\color{black} \newline
Taking into account all of these factors, it is not a surprise that unconsciously we are creating the perfect cocktail for the proliferation of hate speech, insults and personal denigration. The European Commission against Racism and Intolerance General Policy Recommendation No. 15 on combating hate speech states that "the gravity of hate speech targeting women both on account of their sex, gender and/or gender identity and when this is coupled with one or more of their other characteristics". In particular, in March 2018, following the release of its "Toxic Twitter" report\footnote{From \url{https://www.business-humanrights.org/en/latest-news/amnesty-intl-finds-high-levels-of-online-abuse-of-women-on-twitter-incl-company-comments/}, last visited 16/08/24}, Amnesty International initiated Troll Patrol, "a global crowdsourcing effort to demonstrate the scale and nature of abuse that women continue to experience on Twitter". They show evidence that women are more prone to suffer sexism and hate speech in Twitter, as "7.1\% of tweets sent to the women in the study were problematic or abusive. This amounts to 1.1 million problematic or abusive mentions of these 778 women across the year, or one every 30 seconds on average".

Therefore, the motivation of this work is to create reliable and unbiased automatic detection models that help identify hate speech towards women in social networks \textit{(X/Twitter)}. In particular, by means of artificial intelligence (AI) and Natural Language Processing (NLP) algorithms. The early days of AI started back in 1950s-60s where Perceptron Mark I was the first Artificial Neural Network created \color{blue} [C. Giattino, E. Mathieu, J. Broden and M. Roser, in Our World in Data (https://ourworldindata.org/artificial-intelligence, 2022).] \color{black}. But it is not until 2010s with the creation of AlexNet, that this field starts developing at a high and fast pace, thanks to the recent advances in computing power and data availability. The publication of the Transformer architecture by Google \color{blue} [] \color{black}, and recently, with the release of GPT-3 and GPT-4 models, and ChatGPT by OpenAI have led to considerable progress in the field. NLP arised around 2018 and in a few years it was able to perform better than humans, compared with handwriting or image recognition that took longer \color{blue} [C. Giattino, E. Mathieu, J. Broden and M. Roser, in Our World in Data (https://ourworldindata.org/artificial-intelligence, 2022).] \color{black}. \newline \color{red}[definition of NLP]\color{black} \newline
A typical NLP pipeline consists on a series of sequential steps: Text Information, Segmentation and Tokenization, Text Cleaning, Vectorization and Feature Engineering, Text Lemmatization and Steaming, Machine Learning (ML) Algorithms, and Interpretation of the Result \color{blue} [I. Bobriakov, in Data Science Central (https://www.datasciencecentral.com/top-nlp-algorithms-amp-concepts/, 2019)] \color{black}.